\section{Introduction}


\subsection{Motivation}

The rise of decentralized AI computation introduces a fundamental challenge: how to verify that computations executed on remote hardware are trustworthy and tamper-resistant. Traditional cloud computing relies on trust in centralized providers, but decentralized networks require cryptographic proof that code executed inside specific hardware enclaves.

Trusted Execution Environments (TEEs) provide hardware-based isolation and attestation, but existing blockchain systems lack a unified attestation layer capable of:

\begin{itemize}
\item \textbf{Multi-vendor support}: Verifying attestations from Intel SGX, AMD SEV-SNP, NVIDIA H100 CC, and ARM CCA simultaneously
\item \textbf{Compute class diversity}: Supporting CPU-only, GPU-accelerated, and hybrid CPU+GPU workloads
\item \textbf{Single source of truth}: Providing a network-wide attestation registry accessible to all appchains
\item \textbf{Economic separation}: Distinguishing security tokens (LUX) from application tokens (AI, ZOO)
\item \textbf{Performance at scale}: Sub-100ms verification for real-time AI inference workloads
\end{itemize}

A-Chain addresses these challenges by providing a dedicated blockchain layer for attestation verification, serving as the foundation for trustless AI compute across the Lux Network.

\subsection{Key Contributions}

This paper makes the following contributions:

\begin{enumerate}
\item A unified attestation protocol supporting all major TEE vendors and compute classes
\item Gas-efficient on-chain verification using EVM precompiles at address 0xAAA
\item Global attestation registry accessible to all Lux appchains
\item Economic model separating security budget (LUX) from application incentives (AI, ZOO)
\item Performance benchmarks showing sub-100ms attestation verification
\item Integration architecture for AI orchestration layers (Hanzo.Network, Zoo.Network)
\end{enumerate}

